\section{Local forms and the outer reduction}\label{sec:outer}
Fix a dyadic lattice on $\R$.  The depth increases by one when the interval
length is halved.  For an interval $I$ of length $r$, define
\[
 H_I(z_0,z_1,z_2,z_3)
 =r^{-2}\int_{\R^2}w_I(x,t)\prod_{j=0}^3z_j(x+jt)\,dx\,dt.
\]
We fix constants $c_0>0$ and $C_0\ge1$ and call a kernel admissible when
it has the following properties:
\begin{enumerate}
\item Its support consists of points for which all four $x+jt$ lie in
$I$ at distance at least $c_0r$ from its endpoints.
\item For every $u\in\R$ and every $j\in\{0,1,2,3\}$,
\[
 \int_{\R}w_I(u-jt,t)\,dt=0.
\]
\item Every derivative of total order $n$ is bounded in absolute value by
$r^{-n}(C_0(n+2))^{C_0(n+2)}$.
\end{enumerate}
Kernels at different intervals may be chosen independently and may be
multiplied by unimodular scalars.  Constants in the proof may depend on
these fixed admissibility parameters.  For vector-valued functions the
form uses only their physical scalar components.

\begin{lemma}[Local absolute estimate]\label{lem:local-absolute}
If $p_j\ge1$ and $\sum_{j=0}^3p_j^{-1}\le2$, then
\[
 r^{-2}\int |w_I(x,t)|\prod_{j=0}^3\norm{z_j(x+jt)}\,dx\,dt
 \lesssim\prod_{j=0}^3\norm{z_j}_{p_j,I}.
\]
\end{lemma}
\begin{proof}
For distinct indices $i,j$, the map $(x,t)\mapsto(y_i,y_j)$ has
nonzero constant determinant.  The estimate follows immediately when
$p_i=p_j=1$ and the other two exponents are infinity.  The estimate with
all exponents infinity follows from the support and kernel bounds.
The endpoints with just one exponent equal to one follow from the
two-exponent bounds and normalized norm monotonicity.
Interpolation, or log-convexity applied to this positive form, fills the
polytope $0\le p_j^{-1}\le1$, $\sum p_j^{-1}\le2$.
\end{proof}

Let $M$ be the uncentered Hardy--Littlewood maximal operator and put
$M_qf=(M(|f|^q))^{1/q}$.  For $2<q<3$, let $C_N=C_N(q)$ be the least
constant such that
\[
 \sum_{I\in\mathcal I}|I|\,|H_I(f_0,f_1,f_2,f_3)|
 \le C_N\int_{\R}\prod_{j=0}^3 M_qf_j
\]
for every finite collection $\mathcal I$ in at most $N$ consecutive
depths, every lattice, every admissible kernel family, and compactly
supported $L^q$ inputs.  Lemma~\ref{lem:local-absolute}, with all
exponents two, gives $C_N\lesssim N$.  Indeed the local $L^2$ norms
are at most the local $L^q$ norms, whose product is at most
$\prod_jM_qf_j$ at every point of $I$; sum over disjoint intervals at
each depth.

\begin{theorem}[Uniform local estimate]\label{thm:local}
For some $q\in(2,3)$, the constants $C_N(q)$ are bounded independently
of $N$.
\end{theorem}

It is useful to isolate the normalized assertion that will be proved
in Section~\ref{sec:completion}.  A rooted subtree contains all
ancestors between each of its cells and its root; the summation may use
any subset of those cells.

\begin{proposition}[Normalized absorption estimate]\label{prop:normalized}
For a suitable fixed $q\in(2,3)$ and every sufficiently small fixed
$c>0$, there is an $A<\infty$, independent of $N$, with the following
property.  Suppose a finite rooted subtree has root $I_*$, all summation
cells lie in at most $N$ consecutive depths, and
\[
 \norm{f_j}_{q,I}\le1
 \qquad(0\le j\le3, I\text{ in the subtree}).
\]
Then
\[
 \sum_{I\text{ a summation cell}}|I|\,|H_I(f)|
 \le (A+cC_N)|I_*|.
\]
\end{proposition}

The rest of the proof establishes Proposition~\ref{prop:normalized}.
At this stage $C_N$ is finite, with the preliminary bound $C_N\lesssim N$.
The constant $A$ may depend on the fixed choices of $q$, $c$, and the
accuracy parameters, but is independent of $N$.  The stopping argument
below gives $C_N\le C A+C cC_N$; choosing $c$ sufficiently small then
gives the uniform bound.  Thus the long construction needs a small
coefficient of the finite constant $C_N$, which will be absorbed only
after Proposition~\ref{prop:normalized} is proved.

\begin{lemma}[Outer reduction]\label{lem:outer}
Proposition~\ref{prop:normalized} implies Theorem~\ref{thm:local},
and Theorem~\ref{thm:local} implies Theorem~\ref{thm:main}.
\end{lemma}
\begin{proof}
We first remove the normalization.  Start with disjoint roots at the
largest interval size occurring in the finite collection.  At a root
$J$, put $a_{j,J}=\norm{f_j}_{q,J}$.  A zero root size gives zero
contribution below that root.  Otherwise stop at the first strict
descendants $J'$ for which one average of $|f_j|^q$ exceeds
$16a_{j,J}^q$.  For each fixed $j$ the union of its first stops has length
at most $|J|/16$, so the union of all stops has length at most $|J|/4$.
On the intervening tree, division of $f_j$ by $16^{1/q}a_{j,J}$ makes
all required local sizes at most one.  Proposition~\ref{prop:normalized}
therefore bounds this tree by
\[
 16^{4/q}(A+cC_N)|J|\prod_ja_{j,J}.
\]
Repeat at the stopped roots.  The sets
$E_J=J\setminus\bigcup J'$ are disjoint over this recursive family,
have measure at least $3|J|/4$, and satisfy
$M_qf_j(x)\ge a_{j,J}$ for $x\in E_J$.  Hence
\[
 \sum_J |J|\prod_ja_{j,J}
 \le\tfrac43\int\prod_jM_qf_j.
\]
Taking the supremum gives
$C_N\le C A+C cC_N$, where $C$ is bounded uniformly for $2<q<3$.
Choose $c$ so that $Cc<1/2$ and absorb.  This proves
Theorem~\ref{thm:local}.

To recover the continuous kernel, take an odd compactly supported
Gevrey function $\phi$ with $\phi^{(3)}(0)\ne0$, and a compactly
supported Gevrey function $\omega$ such that
$\sum_{m\in\Z}\omega(X-m)=1$.  Set
\[
 W(X,S)=\prod_{j=0}^3(\partial_S-j\partial_X)
                 [\omega(X)\phi(S)].
\]
Each factor is differentiation in the direction of integration on the
corresponding line $X=u-jS$.  Compact support therefore makes all four
one-point marginals of $W$ vanish.  Spatial summation and a change of
variables give
\[
 \sum_{m\in\Z}W(x/a-m,t/a)=\phi^{(4)}(t/a),\qquad
 \int_0^\infty\phi^{(4)}(t/a)\,\frac{da}{a^2}
 =-\frac{\phi^{(3)}(0)}{t}\quad(t\ne0).
\]
For the second equality one first takes $t>0$ and substitutes $s=t/a$;
for $t<0$ use oddness of $\phi^{(4)}$.

Write $a=b2^n$, $1\le b<2$.  For each fixed $b$, a kernel
$a^{-1}W(x/a-m,t/a)$ can be put into the form $r^{-1}w_I(x,t)$ with
$r=Da$ for a fixed sufficiently large dyadic $D$.  To choose $I$, use
a fixed finite family of shifted lattices.  At scale $r$ their
boundaries may be placed at offsets $0,r/3,-r/3$, with alternating signs
between depths to preserve nesting.  The four progression coordinates
on the support occupy an interval of length $O(a)$; for large $D$, one
of these cells contains that interval with the fixed interior margin.
Each chosen cell receives only a bounded number of translates.  The
rescaled derivative and support bounds are fixed, so choose $c_0,C_0$
once to include them.  A bounded number of contributions at one cell
can be estimated separately.

Theorem~\ref{thm:local}, summed over these lattices, gives a uniform
bound for the absolute sum over any finite range of $n$ and the relevant
spatial translates.  Integrate in $b$ with measure $db/b$.  For compactly
supported bounded $f_0$ and Schwartz $f_1,f_2,f_3$, the resulting bound is
\[
 C\int\prod_{j=0}^3M_qf_j
 \le C_q\norm{f_0}_\infty\prod_{j=1}^3\norm{f_j}_3.
\]
Here $M_q$ is bounded on $L^3$ because $q<3$; use H\"older with the
three exponents equal to three.  Compact truncations of the other
inputs justify use of the defining class for $C_N$: at a finite range
of scales and with compact support in slot zero, every relevant
progression coordinate lies in a fixed compact set.

After spatial summation the kernel is odd.  Pair positive and negative
$t$ near zero; the difference of the two products of Schwartz functions
is $O(|t|)$ uniformly on the compact support of $f_0$.  At large $|t|$
the Schwartz bounds are integrable.  The integrated scale kernels have
absolute value at most $C/|t|$.  Dominated convergence therefore permits
the scale range to tend to all positive $a$ and recovers a nonzero
constant times the principal-value four-linear form.  Reflection of
$t$ changes only its sign and gives the operator in
Theorem~\ref{thm:main}.  Duality against compactly supported bounded
$f_0$ proves the inequality in Theorem~\ref{thm:main}.  Density of Schwartz functions,
trilinearity, and completeness of $L^1$ give the stated extension.
\end{proof}
